\section{コード・データ・テスト}\label{app:code}
\begin{itemize}
\item 実装：\texttt{src/rcsknots/refined\_kauffman/}（\texttt{kauffman\_*.py}、\texttt{koornwinder.py}、\texttt{daha\_dn.py}、\texttt{daha\_cn.py}（C$_N$ 型）、\texttt{daha\_fast.py}、\texttt{flint\_rf*.py}）と \texttt{scripts/refined\_kauffman/rk\_*.py}（Hilbert scheme のセル分解は \texttt{rk\_b\_planB\_cells.py}）。
\item 自動テスト：\texttt{tests/refined\_kauffman/}。保存済みの大きな計算結果を判定し直すテストには \texttt{results} の印を付け、既定の実行から外している。
\item 主張とテストの対応は \texttt{docs/refined\_kauffman/paper/claims.md} に置いた。
\item 以下の式は、保存済みのデータとコードから \texttt{scripts/rk\_report\_tex.py} で生成した（手で写していない）。
\end{itemize}

\section{トーラス結び目の表現 [1]}\label{app:torus}
$g^{\mathrm{true}}_{\yone}$（太字の変数）、$c = 1$ 部分の $G$ 展開、$c = 2$ 部分の $\Pi_m$。
$m = 1$（$T(2,3)$、11 項）：$g^{\mathrm{true}}_{[1]}$ の $\ba$ の冪ごとの係数
\begin{align*}
[\ba^{-6}]\colon\ & \bt^{-4} \\
[\ba^{-5}]\colon\ & \bq \bt^{-2} -\bq \bt^{-4} \\
[\ba^{-4}]\colon\ & -\bq^{2} \bt^{-2} -\bt^{-2} -\bq^{-2} \bt^{-4} \\
[\ba^{-3}]\colon\ & -\bq^{-1} \bt^{-2} +\bq^{-1} \bt^{-4} \\
[\ba^{-2}]\colon\ & \bq^{2} +\bt^{-2} +\bq^{-2} \bt^{-2}
\end{align*}

$m = 1$：$c = 1$ 部分の $G$ 展開（$\nu^\beta$ ごと、$W = (q/t)^{1/4}$）
\begin{align*}
[\nu^{-6}]\colon\ & W^{-8} \\
[\nu^{-4}]\colon\ & -G W^{-6} -2 W^{-4} -W^{-8} \\
[\nu^{-2}]\colon\ & G W^{-2} +1 +2 W^{-4}
\end{align*}

\begin{align*}
\Pi_{1} &= 1
\end{align*}

$m = 2$（$T(2,5)$、27 項）：$g^{\mathrm{true}}_{[1]}$ の $\ba$ の冪ごとの係数
\begin{align*}
[\ba^{-10}]\colon\ & \bt^{-8} \\
[\ba^{-9}]\colon\ & \bq \bt^{-6} -\bq \bt^{-8} \\
[\ba^{-8}]\colon\ & \bq^{2} \bt^{-4} -\bq^{2} \bt^{-6} +\bq^{-2} \bt^{-6} -\bq^{-2} \bt^{-8} \\
[\ba^{-7}]\colon\ & \bq^{3} \bt^{-2} -\bq^{3} \bt^{-4} +\bq^{-1} \bt^{-4} -2 \bq^{-1} \bt^{-6} +\bq^{-1} \bt^{-8} \\
[\ba^{-6}]\colon\ & -\bq^{4} \bt^{-2} -\bq^{2} \bt^{-2} -2 \bt^{-4} +\bt^{-6} -\bq^{-2} \bt^{-4} -\bq^{-4} \bt^{-6} \\
[\ba^{-5}]\colon\ & -\bq \bt^{-2} +\bq \bt^{-4} -\bq^{-3} \bt^{-4} +\bq^{-3} \bt^{-6} \\
[\ba^{-4}]\colon\ & \bq^{4} +\bq^{2} \bt^{-2} +\bt^{-2} +\bq^{-2} \bt^{-4} +\bq^{-4} \bt^{-4}
\end{align*}

$m = 2$：$c = 1$ 部分の $G$ 展開（$\nu^\beta$ ごと、$W = (q/t)^{1/4}$）
\begin{align*}
[\nu^{-10}]\colon\ & W^{-16} \\
[\nu^{-8}]\colon\ & G W^{-10} -G W^{-14} +W^{-8} -W^{-16} \\
[\nu^{-6}]\colon\ & -G^{2} W^{-8} -3 G W^{-6} -2 G W^{-10} -2 W^{-4} -3 W^{-8} \\
[\nu^{-4}]\colon\ & G^{2} W^{-4} +2 G W^{-2} +3 G W^{-6} +1 +2 W^{-4} +2 W^{-8}
\end{align*}

\begin{align*}
\Pi_{2} &= \ba^{2} \bq^{4} \bt^{4} +\ba^{2} \bt^{2} +\bq^{2}
\end{align*}

$m = 3$（$T(2,7)$、49 項）：$g^{\mathrm{true}}_{[1]}$ の $\ba$ の冪ごとの係数
\begin{align*}
[\ba^{-14}]\colon\ & \bt^{-12} \\
[\ba^{-13}]\colon\ & \bq \bt^{-10} -\bq \bt^{-12} \\
[\ba^{-12}]\colon\ & \bq^{2} \bt^{-8} -\bq^{2} \bt^{-10} +\bq^{-2} \bt^{-10} -\bq^{-2} \bt^{-12} \\
[\ba^{-11}]\colon\ & \bq^{3} \bt^{-6} -\bq^{3} \bt^{-8} +\bq^{-1} \bt^{-8} -2 \bq^{-1} \bt^{-10} +\bq^{-1} \bt^{-12} \\
[\ba^{-10}]\colon\ & \bq^{4} \bt^{-4} -\bq^{4} \bt^{-6} +\bt^{-6} -2 \bt^{-8} +\bt^{-10} +\bq^{-4} \bt^{-8} \\
& \quad -\bq^{-4} \bt^{-10} \\
[\ba^{-9}]\colon\ & \bq^{5} \bt^{-2} -\bq^{5} \bt^{-4} +\bq \bt^{-4} -2 \bq \bt^{-6} +\bq \bt^{-8} +\bq^{-3} \bt^{-6} \\
& \quad -2 \bq^{-3} \bt^{-8} +\bq^{-3} \bt^{-10} \\
[\ba^{-8}]\colon\ & -\bq^{6} \bt^{-2} -\bq^{4} \bt^{-2} -2 \bq^{2} \bt^{-4} +\bq^{2} \bt^{-6} -\bt^{-4} -2 \bq^{-2} \bt^{-6} \\
& \quad +\bq^{-2} \bt^{-8} -\bq^{-4} \bt^{-6} -\bq^{-6} \bt^{-8} \\
[\ba^{-7}]\colon\ & -\bq^{3} \bt^{-2} +\bq^{3} \bt^{-4} -\bq^{-1} \bt^{-4} +\bq^{-1} \bt^{-6} -\bq^{-5} \bt^{-6} +\bq^{-5} \bt^{-8} \\
[\ba^{-6}]\colon\ & \bq^{6} +\bq^{4} \bt^{-2} +\bq^{2} \bt^{-2} +\bt^{-4} +\bq^{-2} \bt^{-4} +\bq^{-4} \bt^{-6} \\
& \quad +\bq^{-6} \bt^{-6}
\end{align*}

$m = 3$：$c = 1$ 部分の $G$ 展開（$\nu^\beta$ ごと、$W = (q/t)^{1/4}$）
\begin{align*}
[\nu^{-14}]\colon\ & W^{-24} \\
[\nu^{-12}]\colon\ & G W^{-18} -G W^{-22} +W^{-16} -W^{-24} \\
[\nu^{-10}]\colon\ & G^{2} W^{-12} -G^{2} W^{-16} +2 G W^{-10} -2 G W^{-18} +W^{-8} -W^{-16} \\
[\nu^{-8}]\colon\ & -G^{3} W^{-10} -4 G^{2} W^{-8} -3 G^{2} W^{-12} -5 G W^{-6} -7 G W^{-10} -2 G W^{-14} \\
& \quad -2 W^{-4} -3 W^{-8} -2 W^{-12} \\
[\nu^{-6}]\colon\ & G^{3} W^{-6} +3 G^{2} W^{-4} +4 G^{2} W^{-8} +3 G W^{-2} +6 G W^{-6} +5 G W^{-10} \\
& \quad +1 +2 W^{-4} +2 W^{-8} +2 W^{-12}
\end{align*}

\begin{align*}
\Pi_{3} &= \ba^{4} \bq^{8} \bt^{8} +\ba^{4} \bq^{4} \bt^{6} +\ba^{4} \bt^{4} +\ba^{2} \bq^{6} \bt^{4} +\ba^{2} \bq^{2} \bt^{2} +\bq^{4}
\end{align*}


\section{DAHA：一般の $N$ のreduced 不変量}\label{app:daha}
$\Rred_R(t^N) = T_R^{-6}\times$ 鏡像$(Z)/K_R(t^\rho)$（trefoil、一様な規格化）。$q, t$ の冪は半整数のこともある。
$R_{(1)}(U)$（$U^j$ の係数）：
\begin{align*}
[U^{-5}]\colon\ & -q^{-2} t^{5} +q^{-2} t^{4} \\
[U^{-4}]\colon\ & -q^{-1} t^{4} +q^{-2} t^{4} -q^{-2} t^{3} \\
[U^{-3}]\colon\ & q^{-1} t^{3} -q^{-1} t^{2} \\
[U^{-2}]\colon\ & t^{2} +q^{-1} t
\end{align*}

$R_{(2)}(U)$（$U^j$ の係数）：
\begin{align*}
[U^{-10}]\colon\ & q^{-9} t^{10} -q^{-9} t^{9} -q^{-10} t^{9} +q^{-10} t^{8} \\
[U^{-9}]\colon\ & q^{-7} t^{9} -q^{-7} t^{8} -q^{-9} t^{9} +2 q^{-9} t^{8} -q^{-9} t^{7} +q^{-10} t^{8} \\
& \quad -q^{-10} t^{7} \\
[U^{-8}]\colon\ & q^{-5} t^{8} -q^{-6} t^{8} +q^{-6} t^{7} -2 q^{-7} t^{8} +3 q^{-7} t^{7} -q^{-7} t^{6} \\
& \quad +q^{-8} t^{7} -q^{-8} t^{6} -q^{-9} t^{7} +q^{-9} t^{6} \\
[U^{-7}]\colon\ & -q^{-4} t^{7} +q^{-4} t^{6} -2 q^{-5} t^{7} +2 q^{-5} t^{6} -q^{-6} t^{6} +q^{-6} t^{5} \\
& \quad +q^{-7} t^{7} -3 q^{-7} t^{6} +2 q^{-7} t^{5} -q^{-8} t^{6} +q^{-8} t^{5} \\
[U^{-6}]\colon\ & -q^{-2} t^{6} -q^{-3} t^{6} +q^{-4} t^{6} -2 q^{-4} t^{5} +q^{-5} t^{6} -3 q^{-5} t^{5} \\
& \quad +q^{-5} t^{4} +q^{-6} t^{6} -q^{-6} t^{4} +q^{-7} t^{5} -q^{-7} t^{4} \\
[U^{-5}]\colon\ & q^{-2} t^{5} -q^{-2} t^{4} +q^{-3} t^{5} -q^{-3} t^{4} +q^{-4} t^{4} -q^{-4} t^{3} \\
& \quad +q^{-5} t^{4} -q^{-5} t^{3} \\
[U^{-4}]\colon\ & t^{4} +q^{-2} t^{3} +q^{-3} t^{3} +q^{-4} t^{2}
\end{align*}

$R_{(1,1)}(U)$（$U^j$ の係数）：
\begin{align*}
[U^{-10}]\colon\ & q^{-4} t^{16} -q^{-4} t^{15} -q^{-5} t^{15} +q^{-5} t^{14} \\
[U^{-9}]\colon\ & q^{-3} t^{15} -q^{-3} t^{13} -q^{-4} t^{14} +q^{-4} t^{13} +q^{-4} t^{12} -q^{-4} t^{11} \\
[U^{-8}]\colon\ & q^{-2} t^{13} -2 q^{-3} t^{13} +2 q^{-3} t^{11} +q^{-4} t^{12} -q^{-4} t^{11} -q^{-4} t^{10} \\
& \quad +q^{-4} t^{9} \\
[U^{-7}]\colon\ & -q^{-2} t^{12} -q^{-2} t^{11} +q^{-2} t^{10} +q^{-2} t^{9} +q^{-3} t^{11} -q^{-3} t^{10} \\
& \quad -2 q^{-3} t^{9} +q^{-3} t^{8} +q^{-3} t^{7} \\
[U^{-6}]\colon\ & -q^{-1} t^{10} -q^{-1} t^{9} +q^{-2} t^{10} +q^{-2} t^{9} -3 q^{-2} t^{8} -q^{-2} t^{7} \\
& \quad +q^{-3} t^{9} +q^{-3} t^{7} -q^{-3} t^{6} -q^{-3} t^{5} \\
[U^{-5}]\colon\ & q^{-1} t^{8} -q^{-1} t^{6} +q^{-2} t^{6} -q^{-2} t^{4} \\
[U^{-4}]\colon\ & t^{6} +q^{-1} t^{5} +q^{-1} t^{4} +q^{-2} t^{2}
\end{align*}

被覆の随伴 $(1;1)$ の $R_{\mathrm{adj}}$（$A^{1/2} = \nu$ の冪の係数）：
\begin{align*}
[A^{-5}]\colon\ & q^{-4} t^{5} -q^{-4} t^{4} -q^{-5} t^{5} +q^{-5} t^{4} \\
[A^{-9/2}]\colon\ & 0 \\
[A^{-4}]\colon\ & q^{-2} t^{4} -2 q^{-3} t^{4} +2 q^{-3} t^{3} +q^{-4} t^{4} -2 q^{-4} t^{3} +q^{-4} t^{2} \\
[A^{-7/2}]\colon\ & 0 \\
[A^{-3}]\colon\ & -2 q^{-1} t^{3} +q^{-2} t^{3} -3 q^{-2} t^{2} +q^{-3} t^{3} +q^{-3} t^{2} -2 q^{-3} t \\
[A^{-5/2}]\colon\ & 0 \\
[A^{-2}]\colon\ & t^{2} +2 q^{-1} t +q^{-2}
\end{align*}


\section{大きな $N$ の 2 箱のクロスキャップ部分}\label{app:largeN}
trefoil の $\hat g_{\ytwo} + \hat g_{\yoneone}$、$(1+t)(\hat g_{\ytwo} - \hat g_{\yoneone})$、普遍因子で割った商（太字の変数）。
整数性（計算結果 \ref{cr:largeN}）と普遍因子（§\ref{sec:universal}）の確認用で、$\ba$ の冪ごとの符号はそろわない（計算結果 \ref{cr:signs}）。正値性が見えるセクターと見えないセクターの比較は\ref{app:cover}。
$\hat g_{(2)} + \hat g_{(1,1)}$（整数 Laurent、262 項）：
\begin{align*}
[\ba^{-12}]\colon\ & -\bq^{11} \bt^{-9} -\bq^{9} \bt^{-9} -\bq^{7} \bt^{-9} -\bq^{5} \bt^{-9} -\bq^{3} \bt^{-9} -\bq \bt^{-9} \\
& \quad +\bq^{-1} \bt^{-9} +\bq^{-3} \bt^{-11} +\bq^{-5} \bt^{-13} +\bq^{-7} \bt^{-15} +\bq^{-9} \bt^{-17} +\bq^{-11} \bt^{-19} \\
[\ba^{-11}]\colon\ & -\bq^{12} \bt^{-7} +\bq^{12} \bt^{-9} -2 \bq^{10} \bt^{-7} +\bq^{10} \bt^{-9} -2 \bq^{8} \bt^{-7} -2 \bq^{6} \bt^{-7} \\
& \quad -\bq^{6} \bt^{-9} -2 \bq^{4} \bt^{-7} -2 \bq^{4} \bt^{-9} -2 \bq^{2} \bt^{-7} -2 \bq^{2} \bt^{-9} -\bq^{2} \bt^{-11} \\
& \quad -4 \bt^{-9} -2 \bt^{-11} -4 \bq^{-2} \bt^{-11} -\bq^{-2} \bt^{-13} -4 \bq^{-4} \bt^{-13} -3 \bq^{-6} \bt^{-15} \\
& \quad -2 \bq^{-8} \bt^{-17} -\bq^{-10} \bt^{-19} \\
[\ba^{-10}]\colon\ & \bq^{13} \bt^{-7} -\bq^{11} \bt^{-5} +3 \bq^{11} \bt^{-7} -2 \bq^{9} \bt^{-5} +3 \bq^{9} \bt^{-7} +2 \bq^{9} \bt^{-9} \\
& \quad -3 \bq^{7} \bt^{-5} +3 \bq^{7} \bt^{-7} +3 \bq^{7} \bt^{-9} -4 \bq^{5} \bt^{-5} +2 \bq^{5} \bt^{-7} +5 \bq^{5} \bt^{-9} \\
& \quad -4 \bq^{3} \bt^{-5} +\bq^{3} \bt^{-7} +5 \bq^{3} \bt^{-9} +\bq^{3} \bt^{-11} -2 \bq \bt^{-5} -5 \bq \bt^{-7} \\
& \quad +5 \bq \bt^{-9} +3 \bq \bt^{-11} -4 \bq^{-1} \bt^{-7} -2 \bq^{-1} \bt^{-9} +4 \bq^{-1} \bt^{-11} +\bq^{-1} \bt^{-13} \\
& \quad -5 \bq^{-3} \bt^{-9} -2 \bq^{-3} \bt^{-11} +4 \bq^{-3} \bt^{-13} -4 \bq^{-5} \bt^{-11} -\bq^{-5} \bt^{-13} +2 \bq^{-5} \bt^{-15} \\
& \quad -4 \bq^{-7} \bt^{-13} +\bq^{-7} \bt^{-17} -3 \bq^{-9} \bt^{-15} -2 \bq^{-11} \bt^{-17} -\bq^{-13} \bt^{-19} \\
[\ba^{-9}]\colon\ & 2 \bq^{12} \bt^{-5} -\bq^{12} \bt^{-7} +4 \bq^{10} \bt^{-5} -\bq^{10} \bt^{-9} +5 \bq^{8} \bt^{-5} +2 \bq^{8} \bt^{-7} \\
& \quad -\bq^{8} \bt^{-9} +6 \bq^{6} \bt^{-5} +5 \bq^{6} \bt^{-7} -\bq^{6} \bt^{-9} +6 \bq^{4} \bt^{-5} +9 \bq^{4} \bt^{-7} \\
& \quad +4 \bq^{2} \bt^{-5} +13 \bq^{2} \bt^{-7} +4 \bq^{2} \bt^{-9} -\bq^{2} \bt^{-11} +12 \bt^{-7} +10 \bt^{-9} \\
& \quad +15 \bq^{-2} \bt^{-9} +6 \bq^{-2} \bt^{-11} -\bq^{-2} \bt^{-13} +13 \bq^{-4} \bt^{-11} +2 \bq^{-4} \bt^{-13} +10 \bq^{-6} \bt^{-13} \\
& \quad +6 \bq^{-8} \bt^{-15} +3 \bq^{-10} \bt^{-17} +\bq^{-12} \bt^{-19} \\
[\ba^{-8}]\colon\ & -\bq^{13} \bt^{-5} +\bq^{11} \bt^{-3} -2 \bq^{11} \bt^{-5} -\bq^{11} \bt^{-7} +2 \bq^{9} \bt^{-3} -\bq^{9} \bt^{-5} \\
& \quad -4 \bq^{9} \bt^{-7} +4 \bq^{7} \bt^{-3} -6 \bq^{7} \bt^{-7} -\bq^{7} \bt^{-9} +5 \bq^{5} \bt^{-3} +4 \bq^{5} \bt^{-5} \\
& \quad -10 \bq^{5} \bt^{-7} -2 \bq^{5} \bt^{-9} +4 \bq^{3} \bt^{-3} +10 \bq^{3} \bt^{-5} -11 \bq^{3} \bt^{-7} -6 \bq^{3} \bt^{-9} \\
& \quad +17 \bq \bt^{-5} -7 \bq \bt^{-7} -10 \bq \bt^{-9} -\bq \bt^{-11} +5 \bq^{-1} \bt^{-5} +14 \bq^{-1} \bt^{-7} \\
& \quad -13 \bq^{-1} \bt^{-9} -5 \bq^{-1} \bt^{-11} +8 \bq^{-3} \bt^{-7} +7 \bq^{-3} \bt^{-9} -11 \bq^{-3} \bt^{-11} -\bq^{-3} \bt^{-13} \\
& \quad +7 \bq^{-5} \bt^{-9} +3 \bq^{-5} \bt^{-11} -7 \bq^{-5} \bt^{-13} +5 \bq^{-7} \bt^{-11} +\bq^{-7} \bt^{-13} -3 \bq^{-7} \bt^{-15} \\
& \quad +4 \bq^{-9} \bt^{-13} -\bq^{-9} \bt^{-17} +2 \bq^{-11} \bt^{-15} +\bq^{-13} \bt^{-17} \\
[\ba^{-7}]\colon\ & -\bq^{12} \bt^{-3} -2 \bq^{10} \bt^{-3} -2 \bq^{10} \bt^{-5} +\bq^{10} \bt^{-7} -4 \bq^{8} \bt^{-3} -4 \bq^{8} \bt^{-5} \\
& \quad +\bq^{8} \bt^{-7} -6 \bq^{6} \bt^{-3} -8 \bq^{6} \bt^{-5} +\bq^{6} \bt^{-7} -6 \bq^{4} \bt^{-3} -15 \bq^{4} \bt^{-5} \\
& \quad -\bq^{4} \bt^{-7} +\bq^{4} \bt^{-9} -2 \bq^{2} \bt^{-3} -22 \bq^{2} \bt^{-5} -7 \bq^{2} \bt^{-7} +2 \bq^{2} \bt^{-9} \\
& \quad -11 \bt^{-5} -22 \bt^{-7} +\bt^{-11} -19 \bq^{-2} \bt^{-7} -12 \bq^{-2} \bt^{-9} +2 \bq^{-2} \bt^{-11} \\
& \quad -17 \bq^{-4} \bt^{-9} -5 \bq^{-4} \bt^{-11} +\bq^{-4} \bt^{-13} -12 \bq^{-6} \bt^{-11} -\bq^{-6} \bt^{-13} -7 \bq^{-8} \bt^{-13} \\
& \quad -3 \bq^{-10} \bt^{-15} -\bq^{-12} \bt^{-17} \\
[\ba^{-6}]\colon\ & \bq^{11} \bt^{-5} -\bq^{9} \bt^{-3} +2 \bq^{9} \bt^{-5} -\bq^{7} \bt^{-1} -2 \bq^{7} \bt^{-3} +3 \bq^{7} \bt^{-5} \\
& \quad +\bq^{7} \bt^{-7} -\bq^{5} \bt^{-1} -6 \bq^{5} \bt^{-3} +5 \bq^{5} \bt^{-5} +3 \bq^{5} \bt^{-7} -11 \bq^{3} \bt^{-3} \\
& \quad +4 \bq^{3} \bt^{-5} +8 \bq^{3} \bt^{-7} -7 \bq \bt^{-3} -7 \bq \bt^{-5} +13 \bq \bt^{-7} +2 \bq \bt^{-9} \\
& \quad -19 \bq^{-1} \bt^{-5} +11 \bq^{-1} \bt^{-7} +7 \bq^{-1} \bt^{-9} -4 \bq^{-3} \bt^{-5} -11 \bq^{-3} \bt^{-7} +12 \bq^{-3} \bt^{-9} \\
& \quad +2 \bq^{-3} \bt^{-11} -5 \bq^{-5} \bt^{-7} -4 \bq^{-5} \bt^{-9} +8 \bq^{-5} \bt^{-11} -3 \bq^{-7} \bt^{-9} -\bq^{-7} \bt^{-11} \\
& \quad +3 \bq^{-7} \bt^{-13} -2 \bq^{-9} \bt^{-11} +\bq^{-9} \bt^{-15} -\bq^{-11} \bt^{-13} \\
[\ba^{-5}]\colon\ & \bq^{10} \bt^{-3} +\bq^{8} \bt^{-1} +2 \bq^{8} \bt^{-3} +2 \bq^{6} \bt^{-1} +5 \bq^{6} \bt^{-3} +2 \bq^{4} \bt^{-1} \\
& \quad +10 \bq^{4} \bt^{-3} +2 \bq^{4} \bt^{-5} -\bq^{4} \bt^{-7} +12 \bq^{2} \bt^{-3} +9 \bq^{2} \bt^{-5} -2 \bq^{2} \bt^{-7} \\
& \quad +3 \bt^{-3} +21 \bt^{-5} -\bt^{-7} -\bt^{-9} +9 \bq^{-2} \bt^{-5} +12 \bq^{-2} \bt^{-7} \\
& \quad -2 \bq^{-2} \bt^{-9} +10 \bq^{-4} \bt^{-7} +4 \bq^{-4} \bt^{-9} -\bq^{-4} \bt^{-11} +6 \bq^{-6} \bt^{-9} +\bq^{-6} \bt^{-11} \\
& \quad +3 \bq^{-8} \bt^{-11} +\bq^{-10} \bt^{-13} \\
[\ba^{-4}]\colon\ & \bq^{5} \bt^{-1} -\bq^{5} \bt^{-5} +\bq^{3} \bt^{-1} +2 \bq^{3} \bt^{-3} -3 \bq^{3} \bt^{-5} +7 \bq \bt^{-3} \\
& \quad -6 \bq \bt^{-5} -\bq \bt^{-7} +3 \bq^{-1} \bt^{-3} +\bq^{-1} \bt^{-5} -4 \bq^{-1} \bt^{-7} +6 \bq^{-3} \bt^{-5} \\
& \quad -5 \bq^{-3} \bt^{-7} -\bq^{-3} \bt^{-9} +\bq^{-5} \bt^{-5} +2 \bq^{-5} \bt^{-7} -3 \bq^{-5} \bt^{-9} +\bq^{-7} \bt^{-7} \\
& \quad -\bq^{-7} \bt^{-11} \\
[\ba^{-3}]\colon\ & -\bq^{6} \bt^{-1} -2 \bq^{4} \bt^{-1} -\bq^{4} \bt^{-3} -\bq^{2} \bt^{-1} -5 \bq^{2} \bt^{-3} +\bq^{2} \bt^{-5} \\
& \quad -6 \bt^{-3} -\bq^{-2} \bt^{-3} -5 \bq^{-2} \bt^{-5} +\bq^{-2} \bt^{-7} -2 \bq^{-4} \bt^{-5} -\bq^{-4} \bt^{-7} \\
& \quad -\bq^{-6} \bt^{-7} \\
[\ba^{-2}]\colon\ & -\bq^{-1} \bt^{-3} +\bq^{-1} \bt^{-5} \\
[\ba^{-1}]\colon\ & -\bt^{-1} +\bt^{-3}
\end{align*}

$(1 + t)(\hat g_{(2)} - \hat g_{(1,1)})$（整数 Laurent、290 項）：
\begin{align*}
[\ba^{-12}]\colon\ & -\bq^{13} \bt^{-9} -\bq \bt^{-9} -\bq^{-1} \bt^{-9} +\bq^{-1} \bt^{-11} +\bq^{-3} \bt^{-11} -\bq^{-3} \bt^{-13} \\
& \quad -\bq^{-5} \bt^{-13} +\bq^{-5} \bt^{-15} +\bq^{-7} \bt^{-15} -\bq^{-7} \bt^{-17} -\bq^{-9} \bt^{-17} +\bq^{-9} \bt^{-19} \\
& \quad +\bq^{-11} \bt^{-19} \\
[\ba^{-11}]\colon\ & -\bq^{14} \bt^{-7} +\bq^{14} \bt^{-9} -\bq^{12} \bt^{-7} -\bq^{10} \bt^{-9} -\bq^{8} \bt^{-9} -\bq^{6} \bt^{-9} \\
& \quad -\bq^{4} \bt^{-11} -\bq^{2} \bt^{-11} -\bt^{-13} -\bq^{-2} \bt^{-13} -\bq^{-4} \bt^{-15} -\bq^{-6} \bt^{-15} \\
& \quad -\bq^{-8} \bt^{-19} -\bq^{-10} \bt^{-19} \\
[\ba^{-10}]\colon\ & \bq^{15} \bt^{-7} -\bq^{13} \bt^{-5} +2 \bq^{13} \bt^{-7} -\bq^{11} \bt^{-5} +2 \bq^{11} \bt^{-9} -\bq^{9} \bt^{-5} \\
& \quad +\bq^{9} \bt^{-9} -\bq^{7} \bt^{-5} -\bq^{7} \bt^{-7} +2 \bq^{7} \bt^{-9} -\bq^{5} \bt^{-7} +\bq^{5} \bt^{-11} \\
& \quad -\bq^{3} \bt^{-5} -\bq^{3} \bt^{-7} +2 \bq^{3} \bt^{-11} -\bq \bt^{-5} +\bq \bt^{-7} -\bq \bt^{-11} \\
& \quad +\bq \bt^{-13} +\bq^{-1} \bt^{-7} -2 \bq^{-1} \bt^{-9} -2 \bq^{-1} \bt^{-11} +3 \bq^{-1} \bt^{-13} -\bq^{-3} \bt^{-9} \\
& \quad +\bq^{-3} \bt^{-13} -\bq^{-5} \bt^{-13} +\bq^{-5} \bt^{-17} -\bq^{-7} \bt^{-15} +\bq^{-7} \bt^{-17} -\bq^{-9} \bt^{-15} \\
& \quad -\bq^{-11} \bt^{-19} -\bq^{-13} \bt^{-19} \\
[\ba^{-9}]\colon\ & 2 \bq^{14} \bt^{-5} -\bq^{14} \bt^{-7} +2 \bq^{12} \bt^{-5} +\bq^{12} \bt^{-7} -\bq^{12} \bt^{-9} +\bq^{10} \bt^{-5} \\
& \quad +2 \bq^{10} \bt^{-7} +\bq^{8} \bt^{-5} +3 \bq^{8} \bt^{-7} +4 \bq^{6} \bt^{-7} +\bq^{6} \bt^{-9} +2 \bq^{4} \bt^{-7} \\
& \quad +4 \bq^{4} \bt^{-9} -\bq^{4} \bt^{-11} +\bq^{2} \bt^{-7} +4 \bq^{2} \bt^{-9} +\bq^{2} \bt^{-11} +3 \bt^{-9} \\
& \quad +4 \bt^{-11} -\bt^{-13} +\bq^{-2} \bt^{-9} +3 \bq^{-2} \bt^{-11} +\bq^{-2} \bt^{-13} +\bq^{-4} \bt^{-11} \\
& \quad +4 \bq^{-4} \bt^{-13} +2 \bq^{-6} \bt^{-13} +2 \bq^{-6} \bt^{-15} +2 \bq^{-8} \bt^{-15} +\bq^{-8} \bt^{-17} +\bq^{-10} \bt^{-17} \\
& \quad +\bq^{-10} \bt^{-19} +\bq^{-12} \bt^{-19} \\
[\ba^{-8}]\colon\ & -\bq^{15} \bt^{-5} +\bq^{13} \bt^{-3} -\bq^{13} \bt^{-5} -\bq^{13} \bt^{-7} +\bq^{11} \bt^{-3} +\bq^{11} \bt^{-5} \\
& \quad -3 \bq^{11} \bt^{-7} +2 \bq^{9} \bt^{-3} +\bq^{9} \bt^{-5} -2 \bq^{9} \bt^{-7} -\bq^{9} \bt^{-9} +\bq^{7} \bt^{-3} \\
& \quad +4 \bq^{7} \bt^{-5} -4 \bq^{7} \bt^{-7} -\bq^{7} \bt^{-9} +6 \bq^{5} \bt^{-5} -\bq^{5} \bt^{-7} -4 \bq^{5} \bt^{-9} \\
& \quad +4 \bq^{3} \bt^{-5} +\bq^{3} \bt^{-7} -4 \bq^{3} \bt^{-9} -\bq^{3} \bt^{-11} +2 \bq \bt^{-5} +4 \bq \bt^{-7} \\
& \quad -\bq \bt^{-9} -4 \bq \bt^{-11} +\bq^{-1} \bt^{-5} +\bq^{-1} \bt^{-7} +3 \bq^{-1} \bt^{-9} -4 \bq^{-1} \bt^{-11} \\
& \quad -\bq^{-1} \bt^{-13} -\bq^{-3} \bt^{-7} +6 \bq^{-3} \bt^{-9} -4 \bq^{-3} \bt^{-13} +3 \bq^{-5} \bt^{-9} -2 \bq^{-5} \bt^{-13} \\
& \quad -\bq^{-5} \bt^{-15} -\bq^{-7} \bt^{-11} +3 \bq^{-7} \bt^{-13} -\bq^{-7} \bt^{-15} -\bq^{-7} \bt^{-17} +2 \bq^{-9} \bt^{-13} \\
& \quad -\bq^{-9} \bt^{-17} +\bq^{-11} \bt^{-17} +\bq^{-13} \bt^{-17} \\
[\ba^{-7}]\colon\ & -\bq^{14} \bt^{-3} -\bq^{12} \bt^{-3} -2 \bq^{12} \bt^{-5} +\bq^{12} \bt^{-7} -2 \bq^{10} \bt^{-3} -2 \bq^{10} \bt^{-5} \\
& \quad -2 \bq^{8} \bt^{-3} -4 \bq^{8} \bt^{-5} -7 \bq^{6} \bt^{-5} -2 \bq^{6} \bt^{-7} +\bq^{6} \bt^{-9} -5 \bq^{4} \bt^{-5} \\
& \quad -6 \bq^{4} \bt^{-7} +\bq^{4} \bt^{-9} -\bq^{2} \bt^{-5} -9 \bq^{2} \bt^{-7} -2 \bq^{2} \bt^{-9} +\bq^{2} \bt^{-11} \\
& \quad +\bt^{-5} -7 \bt^{-7} -6 \bt^{-9} +\bt^{-11} -\bq^{-2} \bt^{-7} -7 \bq^{-2} \bt^{-9} \\
& \quad -3 \bq^{-2} \bt^{-11} +\bq^{-2} \bt^{-13} -3 \bq^{-4} \bt^{-9} -5 \bq^{-4} \bt^{-11} -2 \bq^{-6} \bt^{-11} -4 \bq^{-6} \bt^{-13} \\
& \quad -3 \bq^{-8} \bt^{-13} -\bq^{-8} \bt^{-15} -\bq^{-10} \bt^{-15} -\bq^{-10} \bt^{-17} -\bq^{-12} \bt^{-17} \\
[\ba^{-6}]\colon\ & \bq^{13} \bt^{-5} -\bq^{11} \bt^{-3} +\bq^{11} \bt^{-5} -\bq^{9} \bt^{-1} -\bq^{9} \bt^{-3} +\bq^{9} \bt^{-5} \\
& \quad +\bq^{9} \bt^{-7} -4 \bq^{7} \bt^{-3} +2 \bq^{7} \bt^{-5} +2 \bq^{7} \bt^{-7} -4 \bq^{5} \bt^{-3} -\bq^{5} \bt^{-5} \\
& \quad +5 \bq^{5} \bt^{-7} -7 \bq^{3} \bt^{-5} +5 \bq^{3} \bt^{-7} +2 \bq^{3} \bt^{-9} +\bq \bt^{-3} -9 \bq \bt^{-5} \\
& \quad +3 \bq \bt^{-7} +5 \bq \bt^{-9} -\bq^{-1} \bt^{-5} -4 \bq^{-1} \bt^{-7} +3 \bq^{-1} \bt^{-9} +2 \bq^{-1} \bt^{-11} \\
& \quad +\bq^{-3} \bt^{-5} -5 \bq^{-3} \bt^{-7} -2 \bq^{-3} \bt^{-9} +6 \bq^{-3} \bt^{-11} -\bq^{-5} \bt^{-7} -3 \bq^{-5} \bt^{-9} \\
& \quad +3 \bq^{-5} \bt^{-11} +\bq^{-5} \bt^{-13} -\bq^{-7} \bt^{-9} -\bq^{-7} \bt^{-11} +\bq^{-7} \bt^{-13} +\bq^{-7} \bt^{-15} \\
& \quad -\bq^{-9} \bt^{-13} +\bq^{-9} \bt^{-15} -\bq^{-11} \bt^{-13} \\
[\ba^{-5}]\colon\ & \bq^{12} \bt^{-3} +\bq^{10} \bt^{-1} +\bq^{10} \bt^{-3} +\bq^{8} \bt^{-1} +3 \bq^{8} \bt^{-3} +5 \bq^{6} \bt^{-3} \\
& \quad +2 \bq^{6} \bt^{-5} -\bq^{6} \bt^{-7} +2 \bq^{4} \bt^{-3} +7 \bq^{4} \bt^{-5} -\bq^{4} \bt^{-7} -\bq^{2} \bt^{-3} \\
& \quad +10 \bq^{2} \bt^{-5} +\bq^{2} \bt^{-7} -\bq^{2} \bt^{-9} -\bt^{-3} +4 \bt^{-5} +7 \bt^{-7} \\
& \quad -\bt^{-9} -\bq^{-2} \bt^{-5} +8 \bq^{-2} \bt^{-7} +2 \bq^{-2} \bt^{-9} -\bq^{-2} \bt^{-11} +2 \bq^{-4} \bt^{-7} \\
& \quad +4 \bq^{-4} \bt^{-9} +2 \bq^{-6} \bt^{-9} +2 \bq^{-6} \bt^{-11} +\bq^{-8} \bt^{-11} +\bq^{-8} \bt^{-13} +\bq^{-10} \bt^{-13} \\
[\ba^{-4}]\colon\ & \bq^{7} \bt^{-1} -\bq^{7} \bt^{-5} -\bq^{5} \bt^{-1} +2 \bq^{5} \bt^{-3} -2 \bq^{5} \bt^{-5} -\bq^{3} \bt^{-1} \\
& \quad +5 \bq^{3} \bt^{-3} -3 \bq^{3} \bt^{-5} -\bq^{3} \bt^{-7} +\bq \bt^{-3} +\bq \bt^{-5} -3 \bq \bt^{-7} \\
& \quad -\bq^{-1} \bt^{-3} +5 \bq^{-1} \bt^{-5} -3 \bq^{-1} \bt^{-7} -\bq^{-1} \bt^{-9} +\bq^{-3} \bt^{-7} -2 \bq^{-3} \bt^{-9} \\
& \quad -\bq^{-5} \bt^{-5} +3 \bq^{-5} \bt^{-7} -\bq^{-5} \bt^{-9} -\bq^{-5} \bt^{-11} +\bq^{-7} \bt^{-7} -\bq^{-7} \bt^{-11} \\
[\ba^{-3}]\colon\ & -\bq^{8} \bt^{-1} -\bq^{6} \bt^{-1} -\bq^{6} \bt^{-3} +\bq^{4} \bt^{-1} -4 \bq^{4} \bt^{-3} +\bq^{4} \bt^{-5} \\
& \quad +\bq^{2} \bt^{-1} -3 \bq^{2} \bt^{-3} -\bq^{2} \bt^{-5} +\bt^{-3} -5 \bt^{-5} +\bt^{-7} \\
& \quad +\bq^{-2} \bt^{-3} -3 \bq^{-2} \bt^{-5} -2 \bq^{-4} \bt^{-7} -\bq^{-6} \bt^{-7} \\
[\ba^{-2}]\colon\ & \bq \bt^{-1} -2 \bq \bt^{-3} +\bq \bt^{-5} -\bq^{-1} \bt^{-3} +\bq^{-1} \bt^{-5} \\
[\ba^{-1}]\colon\ & -\bq^{2} \bt^{-1} +\bq^{2} \bt^{-3} -\bt^{-1} +\bt^{-3} \\
[\ba^{0}]\colon\ & \bq \bt^{-1}
\end{align*}

和の商 $X/((q - t)(\nu^2 - t))$（189 項）：
\begin{align*}
[\ba^{-12}]\colon\ & -\bq^{7} \bt^{-3} -\bq^{7} \bt^{-5} -\bq^{7} \bt^{-7} -\bq^{7} \bt^{-9} -\bq^{5} \bt^{-5} -\bq^{5} \bt^{-7} \\
& \quad -\bq^{5} \bt^{-9} -\bq^{3} \bt^{-7} -\bq^{3} \bt^{-9} -\bq \bt^{-9} +\bq^{-3} \bt^{-11} -\bq^{-5} \bt^{-11} \\
& \quad -\bq^{-5} \bt^{-13} -\bq^{-7} \bt^{-13} -\bq^{-7} \bt^{-15} -\bq^{-9} \bt^{-15} -\bq^{-9} \bt^{-17} -\bq^{-11} \bt^{-17} \\
& \quad -\bq^{-11} \bt^{-19} -\bq^{-13} \bt^{-19} -\bq^{-13} \bt^{-21} -\bq^{-15} \bt^{-21} -\bq^{-15} \bt^{-23} \\
[\ba^{-11}]\colon\ & \bq^{8} \bt^{-9} -\bq^{6} \bt^{-3} -\bq^{6} \bt^{-5} -\bq^{6} \bt^{-7} +\bq^{6} \bt^{-9} -2 \bq^{4} \bt^{-5} \\
& \quad -2 \bq^{4} \bt^{-7} -3 \bq^{2} \bt^{-7} -\bq^{2} \bt^{-9} -2 \bt^{-9} +2 \bq^{-2} \bt^{-9} -\bq^{-2} \bt^{-11} \\
& \quad +4 \bq^{-4} \bt^{-11} +4 \bq^{-6} \bt^{-13} +4 \bq^{-8} \bt^{-15} +3 \bq^{-10} \bt^{-17} +2 \bq^{-12} \bt^{-19} +\bq^{-14} \bt^{-21} \\
[\ba^{-10}]\colon\ & \bq^{9} \bt +\bq^{9} \bt^{-1} +\bq^{9} \bt^{-3} +\bq^{9} \bt^{-5} +\bq^{9} \bt^{-7} +2 \bq^{7} \bt^{-1} \\
& \quad +2 \bq^{7} \bt^{-3} +2 \bq^{7} \bt^{-5} +3 \bq^{7} \bt^{-7} +2 \bq^{5} \bt^{-3} +2 \bq^{5} \bt^{-5} +4 \bq^{5} \bt^{-7} \\
& \quad +\bq^{5} \bt^{-9} +2 \bq^{3} \bt^{-5} +5 \bq^{3} \bt^{-7} +2 \bq^{3} \bt^{-9} +6 \bq \bt^{-7} +4 \bq \bt^{-9} \\
& \quad +4 \bq^{-1} \bt^{-7} +5 \bq^{-1} \bt^{-9} +\bq^{-1} \bt^{-11} +2 \bq^{-3} \bt^{-7} +7 \bq^{-3} \bt^{-9} +3 \bq^{-3} \bt^{-11} \\
& \quad +4 \bq^{-5} \bt^{-9} +7 \bq^{-5} \bt^{-11} +\bq^{-5} \bt^{-13} +4 \bq^{-7} \bt^{-11} +6 \bq^{-7} \bt^{-13} +3 \bq^{-9} \bt^{-13} \\
& \quad +4 \bq^{-9} \bt^{-15} +3 \bq^{-11} \bt^{-15} +3 \bq^{-11} \bt^{-17} +2 \bq^{-13} \bt^{-17} +2 \bq^{-13} \bt^{-19} +\bq^{-15} \bt^{-19} \\
& \quad +\bq^{-15} \bt^{-21} \\
[\ba^{-9}]\colon\ & \bq^{8} \bt +\bq^{8} \bt^{-1} +\bq^{8} \bt^{-3} +\bq^{8} \bt^{-5} -\bq^{8} \bt^{-7} +3 \bq^{6} \bt^{-1} \\
& \quad +3 \bq^{6} \bt^{-3} +3 \bq^{6} \bt^{-5} -\bq^{6} \bt^{-7} +5 \bq^{4} \bt^{-3} +5 \bq^{4} \bt^{-5} +8 \bq^{2} \bt^{-5} \\
& \quad +2 \bq^{2} \bt^{-7} -\bq^{2} \bt^{-9} +6 \bt^{-7} -\bt^{-9} -4 \bq^{-2} \bt^{-7} +\bq^{-2} \bt^{-9} \\
& \quad -\bq^{-2} \bt^{-11} -10 \bq^{-4} \bt^{-9} -\bq^{-4} \bt^{-11} -11 \bq^{-6} \bt^{-11} -\bq^{-6} \bt^{-13} -9 \bq^{-8} \bt^{-13} \\
& \quad -6 \bq^{-10} \bt^{-15} -3 \bq^{-12} \bt^{-17} -\bq^{-14} \bt^{-19} \\
[\ba^{-8}]\colon\ & -\bq^{9} \bt^{3} -\bq^{9} \bt -\bq^{9} \bt^{-1} -\bq^{9} \bt^{-3} -\bq^{9} \bt^{-5} -\bq^{7} \bt \\
& \quad -\bq^{7} \bt^{-1} -\bq^{7} \bt^{-3} -2 \bq^{7} \bt^{-5} -\bq^{5} \bt^{-1} -\bq^{5} \bt^{-3} -3 \bq^{5} \bt^{-5} \\
& \quad -\bq^{5} \bt^{-7} -\bq^{3} \bt^{-3} -5 \bq^{3} \bt^{-5} -3 \bq^{3} \bt^{-7} -6 \bq \bt^{-5} -7 \bq \bt^{-7} \\
& \quad -4 \bq^{-1} \bt^{-5} -11 \bq^{-1} \bt^{-7} -2 \bq^{-1} \bt^{-9} -13 \bq^{-3} \bt^{-7} -6 \bq^{-3} \bt^{-9} -3 \bq^{-5} \bt^{-7} \\
& \quad -12 \bq^{-5} \bt^{-9} -2 \bq^{-5} \bt^{-11} -4 \bq^{-7} \bt^{-9} -8 \bq^{-7} \bt^{-11} -3 \bq^{-9} \bt^{-11} -4 \bq^{-9} \bt^{-13} \\
& \quad -2 \bq^{-11} \bt^{-13} -2 \bq^{-11} \bt^{-15} -\bq^{-13} \bt^{-15} -\bq^{-13} \bt^{-17} \\
[\ba^{-7}]\colon\ & -\bq^{8} \bt^{3} -\bq^{8} \bt -\bq^{8} \bt^{-1} -\bq^{8} \bt^{-3} -2 \bq^{6} \bt -2 \bq^{6} \bt^{-1} \\
& \quad -2 \bq^{6} \bt^{-3} -4 \bq^{4} \bt^{-1} -4 \bq^{4} \bt^{-3} -8 \bq^{2} \bt^{-3} -2 \bq^{2} \bt^{-5} +\bq^{2} \bt^{-7} \\
& \quad -7 \bt^{-5} +2 \bt^{-7} +2 \bq^{-2} \bt^{-5} +\bq^{-2} \bt^{-7} +\bq^{-2} \bt^{-9} +7 \bq^{-4} \bt^{-7} \\
& \quad +2 \bq^{-4} \bt^{-9} +9 \bq^{-6} \bt^{-9} +\bq^{-6} \bt^{-11} +6 \bq^{-8} \bt^{-11} +3 \bq^{-10} \bt^{-13} +\bq^{-12} \bt^{-15} \\
[\ba^{-6}]\colon\ & \bq^{3} \bt^{-3} +\bq^{3} \bt^{-5} +\bq \bt^{-3} +3 \bq \bt^{-5} +6 \bq^{-1} \bt^{-5} +\bq^{-1} \bt^{-7} \\
& \quad +3 \bq^{-3} \bt^{-5} +4 \bq^{-3} \bt^{-7} +6 \bq^{-5} \bt^{-7} +\bq^{-5} \bt^{-9} +\bq^{-7} \bt^{-7} +3 \bq^{-7} \bt^{-9} \\
& \quad +\bq^{-9} \bt^{-9} +\bq^{-9} \bt^{-11} \\
[\ba^{-5}]\colon\ & \bq^{4} \bt +\bq^{4} \bt^{-1} +3 \bq^{2} \bt^{-1} +\bq^{2} \bt^{-3} +3 \bt^{-3} -\bt^{-5} \\
& \quad -\bq^{-4} \bt^{-5} -\bq^{-4} \bt^{-7} -2 \bq^{-6} \bt^{-7} -\bq^{-8} \bt^{-9} \\
[\ba^{-4}]\colon\ & -\bq^{-3} \bt^{-5} \\
[\ba^{-3}]\colon\ & -\bq^{-2} \bt^{-3}
\end{align*}

$\hat g_{(1)}$ の商：
\begin{align*}
X_{(1)}/((q - t)(\nu^2 - t)) &= -\ba^{-4} \bq^{-4} \bt^{-4} +\ba^{-5} \bq^{-3} \bt^{-4} +\ba^{-6} \bq^{-4} \bt^{-6}
\end{align*}


\section{被覆の brane--brane と brane--image のセクター}\label{app:cover}
§\ref{sec:split} と §\ref{sec:kn} の根拠。trefoil と $T(2,5)$ の 2 箱について、被覆を brane（$Y$）と image（$Y'$）に分けた生成関数の各セクターを、\cite{KN17} の t 側 $(t^{-1/2} - t^{1/2})\hat f$ で $G$ 展開したもの（$\lambda^\beta$ ごと、$\qr = (q/t)^{1/2}$）。brane--brane は $\hat f_{\ytwo}$, $\hat f_{\yoneone}$（$M \otimes M$ の hatted。一様な規格化の因子を除けば \cite{KN17} の $\hat f$ の鏡像に $R \leftrightarrow R^T$ で一致する、§\ref{sec:split}）、brane--image は $x = W_{\mathrm{adj}} + \kappa - W_{\yone}^2$ を $z_q$ で割ったもの。ブロックの符号の判定は次の表のとおりで、brane--brane は全ブロックでそろい、brane--image は「そろう」「$(\bt^2 - 1)\times$ そろう」「混在」が混じる。
\begin{center}
\footnotesize
\begin{tabular}{llp{9.2cm}}
\toprule
結び目 & セクター & $\lambda^\beta$ ごとの符号 \\
\midrule
$T(2,3)$ & brane--brane $(t^{-1/2} - t^{1/2})\hat f_{\ytwo}$ & $\lambda^{-10}$: そろう、$\lambda^{-8}$: そろう、$\lambda^{-6}$: そろう、$\lambda^{-4}$: そろう、$\lambda^{-2}$: そろう \\
$T(2,3)$ & brane--brane $(t^{-1/2} - t^{1/2})\hat f_{\yoneone}$ & $\lambda^{-10}$: そろう、$\lambda^{-8}$: そろう、$\lambda^{-6}$: そろう、$\lambda^{-4}$: そろう、$\lambda^{-2}$: そろう \\
$T(2,3)$ & brane--image $(t^{-1/2} - t^{1/2})\,x/z_q$ & $\lambda^{-12}$: そろう、$\lambda^{-10}$: $(\bt^2 - 1)\times$ そろう、$\lambda^{-8}$: \textbf{混在}、$\lambda^{-6}$: そろう、$\lambda^{-4}$: そろう、$\lambda^{-2}$: $(\bt^2 - 1)\times$ そろう、$\lambda^{0}$: そろう \\
$T(2,5)$ & brane--brane $(t^{-1/2} - t^{1/2})\hat f_{\ytwo}$ & $\lambda^{-14}$: そろう、$\lambda^{-12}$: そろう、$\lambda^{-10}$: そろう、$\lambda^{-8}$: そろう、$\lambda^{-6}$: そろう \\
$T(2,5)$ & brane--brane $(t^{-1/2} - t^{1/2})\hat f_{\yoneone}$ & $\lambda^{-14}$: そろう、$\lambda^{-12}$: そろう、$\lambda^{-10}$: そろう、$\lambda^{-8}$: そろう、$\lambda^{-6}$: そろう \\
$T(2,5)$ & brane--image $(t^{-1/2} - t^{1/2})\,x/z_q$ & $\lambda^{-20}$: そろう、$\lambda^{-18}$: $(\bt^2 - 1)\times$ そろう、$\lambda^{-16}$: \textbf{混在}、$\lambda^{-14}$: \textbf{混在}、$\lambda^{-12}$: \textbf{混在}、$\lambda^{-10}$: \textbf{混在}、$\lambda^{-8}$: そろう、$\lambda^{-6}$: $(\bt^2 - 1)\times$ そろう、$\lambda^{-4}$: $(\bt^2 - 1)\times$ そろう、$\lambda^{-2}$: $(\bt^2 - 1)\times$ そろう、$\lambda^{0}$: そろう \\
\bottomrule
\end{tabular}
\end{center}

$T(2,3)$、brane--brane $(t^{-1/2} - t^{1/2})\hat f_{\ytwo}$ の $G$ 展開（$\lambda^\beta$ ごと、46 項）
\begin{align*}
[\lambda^{-10}]\colon\ & -G^{2} \qr^{-7} -2 G \qr^{-6} -2 G \qr^{-8} -2 \qr^{-5} -\qr^{-7} -\qr^{-9} \\
[\lambda^{-8}]\colon\ & G^{3} \qr^{-6} +4 G^{2} \qr^{-5} +4 G^{2} \qr^{-7} +6 G \qr^{-4} +9 G \qr^{-6} +5 G \qr^{-8} \\
& \quad +3 \qr^{-3} +7 \qr^{-5} +4 \qr^{-7} +2 \qr^{-9} \\
[\lambda^{-6}]\colon\ & -G^{3} \qr^{-4} -G^{3} \qr^{-6} -3 G^{2} \qr^{-3} -8 G^{2} \qr^{-5} -3 G^{2} \qr^{-7} -3 G \qr^{-2} \\
& \quad -14 G \qr^{-4} -12 G \qr^{-6} -3 G \qr^{-8} -\qr^{-1} -8 \qr^{-3} -9 \qr^{-5} \\
& \quad -5 \qr^{-7} -\qr^{-9} \\
[\lambda^{-4}]\colon\ & G^{3} \qr^{-4} +4 G^{2} \qr^{-3} +4 G^{2} \qr^{-5} +5 G \qr^{-2} +10 G \qr^{-4} +5 G \qr^{-6} \\
& \quad +2 \qr^{-1} +7 \qr^{-3} +5 \qr^{-5} +2 \qr^{-7} \\
[\lambda^{-2}]\colon\ & -G^{2} \qr^{-3} -2 G \qr^{-2} -2 G \qr^{-4} -\qr^{-1} -2 \qr^{-3} -\qr^{-5}
\end{align*}

$T(2,3)$、brane--brane $(t^{-1/2} - t^{1/2})\hat f_{\yoneone}$ の $G$ 展開（$\lambda^\beta$ ごと、27 項）
\begin{align*}
[\lambda^{-10}]\colon\ & -G \qr^{-6} -\qr^{-5} -\qr^{-7} \\
[\lambda^{-8}]\colon\ & G^{2} \qr^{-5} +3 G \qr^{-4} +3 G \qr^{-6} +2 \qr^{-3} +4 \qr^{-5} +2 \qr^{-7} \\
[\lambda^{-6}]\colon\ & -G^{2} \qr^{-3} -G^{2} \qr^{-5} -2 G \qr^{-2} -6 G \qr^{-4} -2 G \qr^{-6} -\qr^{-1} \\
& \quad -5 \qr^{-3} -5 \qr^{-5} -\qr^{-7} \\
[\lambda^{-4}]\colon\ & G^{2} \qr^{-3} +3 G \qr^{-2} +3 G \qr^{-4} +2 \qr^{-1} +4 \qr^{-3} +2 \qr^{-5} \\
[\lambda^{-2}]\colon\ & -G \qr^{-2} -\qr^{-1} -\qr^{-3}
\end{align*}

$T(2,3)$、brane--image $(t^{-1/2} - t^{1/2})\,x/z_q$ の $G$ 展開（$\lambda^\beta$ ごと、29 項）
\begin{align*}
[\lambda^{-12}]\colon\ & -\qr^{-9} \\
[\lambda^{-10}]\colon\ & -G \qr^{-6} +G \qr^{-8} -2 \qr^{-5} +2 \qr^{-9} \\
[\lambda^{-8}]\colon\ & G^{2} \qr^{-5} +4 G \qr^{-4} +3 G \qr^{-6} -G \qr^{-8} +3 \qr^{-3} +7 \qr^{-5} \\
& \quad -\qr^{-9} \\
[\lambda^{-6}]\colon\ & -G^{2} \qr^{-3} -G^{2} \qr^{-5} -2 G \qr^{-2} -8 G \qr^{-4} -2 G \qr^{-6} -\qr^{-1} \\
& \quad -7 \qr^{-3} -8 \qr^{-5} \\
[\lambda^{-4}]\colon\ & G^{2} \qr^{-3} +2 G \qr^{-2} +4 G \qr^{-4} +\qr^{-1} +5 \qr^{-3} +3 \qr^{-5} \\
[\lambda^{-2}]\colon\ & \qr^{-1} -\qr^{-3} \\
[\lambda^{0}]\colon\ & -\qr^{-1}
\end{align*}

$T(2,3)$、同じものを $\lambda^\beta$ ごとに因数分解した形（$\qr^2 - 1 = \bt^2 - 1$）
\begin{align*}
[\lambda^{-12}]\colon\ & \qr^{-9}\,\bigl(-1\bigr) \\
[\lambda^{-10}]\colon\ & \qr^{-9}\,\bigl(- (\qr^2 - 1)(G \qr + 2 \qr^{2} + 2)\bigr) \\
[\lambda^{-8}]\colon\ & \text{既約（上の形のまま）} \\
[\lambda^{-6}]\colon\ & \text{既約（上の形のまま）} \\
[\lambda^{-4}]\colon\ & \text{既約（上の形のまま）} \\
[\lambda^{-2}]\colon\ & \qr^{-3}\,(\qr^2 - 1) \\
[\lambda^{0}]\colon\ & \qr^{-1}\,\bigl(-1\bigr)
\end{align*}

$T(2,5)$、brane--brane $(t^{-1/2} - t^{1/2})\hat f_{\ytwo}$ の $G$ 展開（$\lambda^\beta$ ごと、172 項）
\begin{align*}
[\lambda^{-14}]\colon\ & -G^{6} \qr^{-11} -6 G^{5} \qr^{-10} -6 G^{5} \qr^{-12} -16 G^{4} \qr^{-9} -26 G^{4} \qr^{-11} -15 G^{4} \qr^{-13} \\
& \quad -24 G^{3} \qr^{-8} -49 G^{3} \qr^{-10} -44 G^{3} \qr^{-12} -20 G^{3} \qr^{-14} -22 G^{2} \qr^{-7} -49 G^{2} \qr^{-9} \\
& \quad -54 G^{2} \qr^{-11} -36 G^{2} \qr^{-13} -15 G^{2} \qr^{-15} -12 G \qr^{-6} -28 G \qr^{-8} -30 G \qr^{-10} \\
& \quad -25 G \qr^{-12} -14 G \qr^{-14} -6 G \qr^{-16} -3 \qr^{-5} -7 \qr^{-7} -8 \qr^{-9} \\
& \quad -5 \qr^{-11} -4 \qr^{-13} -2 \qr^{-15} -\qr^{-17} \\
[\lambda^{-12}]\colon\ & G^{7} \qr^{-10} +8 G^{6} \qr^{-9} +8 G^{6} \qr^{-11} +28 G^{5} \qr^{-8} +50 G^{5} \qr^{-10} +27 G^{5} \qr^{-12} \\
& \quad +56 G^{4} \qr^{-7} +133 G^{4} \qr^{-9} +126 G^{4} \qr^{-11} +50 G^{4} \qr^{-13} +70 G^{3} \qr^{-6} +195 G^{3} \qr^{-8} \\
& \quad +240 G^{3} \qr^{-10} +164 G^{3} \qr^{-12} +55 G^{3} \qr^{-14} +55 G^{2} \qr^{-5} +170 G^{2} \qr^{-7} +237 G^{2} \qr^{-9} \\
& \quad +205 G^{2} \qr^{-11} +116 G^{2} \qr^{-13} +36 G^{2} \qr^{-15} +25 G \qr^{-4} +84 G \qr^{-6} +127 G \qr^{-8} \\
& \quad +117 G \qr^{-10} +82 G \qr^{-12} +42 G \qr^{-14} +13 G \qr^{-16} +5 \qr^{-3} +18 \qr^{-5} \\
& \quad +30 \qr^{-7} +28 \qr^{-9} +19 \qr^{-11} +12 \qr^{-13} +6 \qr^{-15} +2 \qr^{-17} \\
[\lambda^{-10}]\colon\ & -G^{7} \qr^{-8} -G^{7} \qr^{-10} -7 G^{6} \qr^{-7} -16 G^{6} \qr^{-9} -7 G^{6} \qr^{-11} -22 G^{5} \qr^{-6} \\
& \quad -72 G^{5} \qr^{-8} -71 G^{5} \qr^{-10} -21 G^{5} \qr^{-12} -40 G^{4} \qr^{-5} -163 G^{4} \qr^{-7} -228 G^{4} \qr^{-9} \\
& \quad -150 G^{4} \qr^{-11} -35 G^{4} \qr^{-13} -45 G^{3} \qr^{-4} -216 G^{3} \qr^{-6} -368 G^{3} \qr^{-8} -335 G^{3} \qr^{-10} \\
& \quad -175 G^{3} \qr^{-12} -35 G^{3} \qr^{-14} -31 G^{2} \qr^{-3} -172 G^{2} \qr^{-5} -334 G^{2} \qr^{-7} -360 G^{2} \qr^{-9} \\
& \quad -252 G^{2} \qr^{-11} -116 G^{2} \qr^{-13} -21 G^{2} \qr^{-15} -12 G \qr^{-2} -77 G \qr^{-4} -166 G \qr^{-6} \\
& \quad -196 G \qr^{-8} -158 G \qr^{-10} -93 G \qr^{-12} -41 G \qr^{-14} -7 G \qr^{-16} -2 \qr^{-1} \\
& \quad -15 \qr^{-3} -35 \qr^{-5} -46 \qr^{-7} -37 \qr^{-9} -25 \qr^{-11} -13 \qr^{-13} \\
& \quad -6 \qr^{-15} -\qr^{-17} \\
[\lambda^{-8}]\colon\ & G^{7} \qr^{-8} +8 G^{6} \qr^{-7} +8 G^{6} \qr^{-9} +28 G^{5} \qr^{-6} +50 G^{5} \qr^{-8} +27 G^{5} \qr^{-10} \\
& \quad +56 G^{4} \qr^{-5} +133 G^{4} \qr^{-7} +126 G^{4} \qr^{-9} +50 G^{4} \qr^{-11} +69 G^{3} \qr^{-4} +195 G^{3} \qr^{-6} \\
& \quad +241 G^{3} \qr^{-8} +164 G^{3} \qr^{-10} +55 G^{3} \qr^{-12} +52 G^{2} \qr^{-3} +166 G^{2} \qr^{-5} +241 G^{2} \qr^{-7} \\
& \quad +208 G^{2} \qr^{-9} +116 G^{2} \qr^{-11} +36 G^{2} \qr^{-13} +22 G \qr^{-2} +77 G \qr^{-4} +127 G \qr^{-6} \\
& \quad +124 G \qr^{-8} +85 G \qr^{-10} +42 G \qr^{-12} +13 G \qr^{-14} +4 \qr^{-1} +15 \qr^{-3} \\
& \quad +28 \qr^{-5} +30 \qr^{-7} +22 \qr^{-9} +13 \qr^{-11} +6 \qr^{-13} +2 \qr^{-15} \\
[\lambda^{-6}]\colon\ & -G^{6} \qr^{-7} -6 G^{5} \qr^{-6} -6 G^{5} \qr^{-8} -16 G^{4} \qr^{-5} -26 G^{4} \qr^{-7} -15 G^{4} \qr^{-9} \\
& \quad -24 G^{3} \qr^{-4} -49 G^{3} \qr^{-6} -44 G^{3} \qr^{-8} -20 G^{3} \qr^{-10} -21 G^{2} \qr^{-3} -49 G^{2} \qr^{-5} \\
& \quad -55 G^{2} \qr^{-7} -36 G^{2} \qr^{-9} -15 G^{2} \qr^{-11} -10 G \qr^{-2} -25 G \qr^{-4} -33 G \qr^{-6} \\
& \quad -27 G \qr^{-8} -14 G \qr^{-10} -6 G \qr^{-12} -2 \qr^{-1} -5 \qr^{-3} -8 \qr^{-5} \\
& \quad -7 \qr^{-7} -5 \qr^{-9} -2 \qr^{-11} -\qr^{-13}
\end{align*}

$T(2,5)$、brane--brane $(t^{-1/2} - t^{1/2})\hat f_{\yoneone}$ の $G$ 展開（$\lambda^\beta$ ごと、133 項）
\begin{align*}
[\lambda^{-14}]\colon\ & -G^{5} \qr^{-10} -5 G^{4} \qr^{-9} -5 G^{4} \qr^{-11} -11 G^{3} \qr^{-8} -17 G^{3} \qr^{-10} -10 G^{3} \qr^{-12} \\
& \quad -13 G^{2} \qr^{-7} -25 G^{2} \qr^{-9} -21 G^{2} \qr^{-11} -10 G^{2} \qr^{-13} -8 G \qr^{-6} -18 G \qr^{-8} \\
& \quad -18 G \qr^{-10} -11 G \qr^{-12} -5 G \qr^{-14} -2 \qr^{-5} -6 \qr^{-7} -5 \qr^{-9} \\
& \quad -4 \qr^{-11} -2 \qr^{-13} -\qr^{-15} \\
[\lambda^{-12}]\colon\ & G^{6} \qr^{-9} +7 G^{5} \qr^{-8} +7 G^{5} \qr^{-10} +21 G^{4} \qr^{-7} +37 G^{4} \qr^{-9} +20 G^{4} \qr^{-11} \\
& \quad +35 G^{3} \qr^{-6} +81 G^{3} \qr^{-8} +75 G^{3} \qr^{-10} +30 G^{3} \qr^{-12} +34 G^{2} \qr^{-5} +94 G^{2} \qr^{-7} \\
& \quad +110 G^{2} \qr^{-9} +73 G^{2} \qr^{-11} +25 G^{2} \qr^{-13} +18 G \qr^{-4} +57 G \qr^{-6} +78 G \qr^{-8} \\
& \quad +62 G \qr^{-10} +34 G \qr^{-12} +11 G \qr^{-14} +4 \qr^{-3} +14 \qr^{-5} +23 \qr^{-7} \\
& \quad +19 \qr^{-9} +12 \qr^{-11} +6 \qr^{-13} +2 \qr^{-15} \\
[\lambda^{-10}]\colon\ & -G^{6} \qr^{-7} -G^{6} \qr^{-9} -6 G^{5} \qr^{-6} -14 G^{5} \qr^{-8} -6 G^{5} \qr^{-10} -16 G^{4} \qr^{-5} \\
& \quad -53 G^{4} \qr^{-7} -52 G^{4} \qr^{-9} -15 G^{4} \qr^{-11} -24 G^{3} \qr^{-4} -99 G^{3} \qr^{-6} -135 G^{3} \qr^{-8} \\
& \quad -88 G^{3} \qr^{-10} -20 G^{3} \qr^{-12} -21 G^{2} \qr^{-3} -102 G^{2} \qr^{-5} -171 G^{2} \qr^{-7} -148 G^{2} \qr^{-9} \\
& \quad -77 G^{2} \qr^{-11} -15 G^{2} \qr^{-13} -10 G \qr^{-2} -55 G \qr^{-4} -109 G \qr^{-6} -113 G \qr^{-8} \\
& \quad -73 G \qr^{-10} -34 G \qr^{-12} -6 G \qr^{-14} -2 \qr^{-1} -12 \qr^{-3} -28 \qr^{-5} \\
& \quad -33 \qr^{-7} -25 \qr^{-9} -13 \qr^{-11} -6 \qr^{-13} -\qr^{-15} \\
[\lambda^{-8}]\colon\ & G^{6} \qr^{-7} +7 G^{5} \qr^{-6} +7 G^{5} \qr^{-8} +21 G^{4} \qr^{-5} +37 G^{4} \qr^{-7} +20 G^{4} \qr^{-9} \\
& \quad +35 G^{3} \qr^{-4} +81 G^{3} \qr^{-6} +75 G^{3} \qr^{-8} +30 G^{3} \qr^{-10} +34 G^{2} \qr^{-3} +93 G^{2} \qr^{-5} \\
& \quad +111 G^{2} \qr^{-7} +73 G^{2} \qr^{-9} +25 G^{2} \qr^{-11} +18 G \qr^{-2} +54 G \qr^{-4} +79 G \qr^{-6} \\
& \quad +64 G \qr^{-8} +34 G \qr^{-10} +11 G \qr^{-12} +4 \qr^{-1} +12 \qr^{-3} +22 \qr^{-5} \\
& \quad +21 \qr^{-7} +13 \qr^{-9} +6 \qr^{-11} +2 \qr^{-13} \\
[\lambda^{-6}]\colon\ & -G^{5} \qr^{-6} -5 G^{4} \qr^{-5} -5 G^{4} \qr^{-7} -11 G^{3} \qr^{-4} -17 G^{3} \qr^{-6} -10 G^{3} \qr^{-8} \\
& \quad -13 G^{2} \qr^{-3} -25 G^{2} \qr^{-5} -21 G^{2} \qr^{-7} -10 G^{2} \qr^{-9} -8 G \qr^{-2} -17 G \qr^{-4} \\
& \quad -19 G \qr^{-6} -11 G \qr^{-8} -5 G \qr^{-10} -2 \qr^{-1} -4 \qr^{-3} -6 \qr^{-5} \\
& \quad -5 \qr^{-7} -2 \qr^{-9} -\qr^{-11}
\end{align*}

$T(2,5)$、brane--image $(t^{-1/2} - t^{1/2})\,x/z_q$ の $G$ 展開（$\lambda^\beta$ ごと、93 項）
\begin{align*}
[\lambda^{-20}]\colon\ & -\qr^{-17} \\
[\lambda^{-18}]\colon\ & -G \qr^{-14} +G \qr^{-16} -\qr^{-13} -\qr^{-15} +2 \qr^{-17} \\
[\lambda^{-16}]\colon\ & -G^{2} \qr^{-11} +G^{2} \qr^{-13} -2 G \qr^{-10} -G \qr^{-12} +4 G \qr^{-14} -G \qr^{-16} \\
& \quad -\qr^{-9} -2 \qr^{-11} +2 \qr^{-13} +2 \qr^{-15} -\qr^{-17} \\
[\lambda^{-14}]\colon\ & -G^{3} \qr^{-8} +G^{3} \qr^{-10} -4 G^{2} \qr^{-7} +5 G^{2} \qr^{-11} -G^{2} \qr^{-13} -5 G \qr^{-6} \\
& \quad -5 G \qr^{-8} +8 G \qr^{-10} +5 G \qr^{-12} -3 G \qr^{-14} -2 \qr^{-5} -4 \qr^{-7} \\
& \quad +\qr^{-9} +7 \qr^{-11} -\qr^{-13} -\qr^{-15} \\
[\lambda^{-12}]\colon\ & G^{4} \qr^{-7} +6 G^{3} \qr^{-6} +5 G^{3} \qr^{-8} -G^{3} \qr^{-10} +12 G^{2} \qr^{-5} +23 G^{2} \qr^{-7} \\
& \quad +4 G^{2} \qr^{-9} -4 G^{2} \qr^{-11} +10 G \qr^{-4} +29 G \qr^{-6} +23 G \qr^{-8} -8 G \qr^{-10} \\
& \quad -4 G \qr^{-12} +3 \qr^{-3} +11 \qr^{-5} +15 \qr^{-7} +4 \qr^{-9} -8 \qr^{-11} \\
[\lambda^{-10}]\colon\ & -G^{4} \qr^{-5} -G^{4} \qr^{-7} -4 G^{3} \qr^{-4} -12 G^{3} \qr^{-6} -4 G^{3} \qr^{-8} -6 G^{2} \qr^{-3} \\
& \quad -28 G^{2} \qr^{-5} -32 G^{2} \qr^{-7} -4 G^{2} \qr^{-9} -4 G \qr^{-2} -24 G \qr^{-4} -44 G \qr^{-6} \\
& \quad -30 G \qr^{-8} +2 G \qr^{-10} -\qr^{-1} -7 \qr^{-3} -17 \qr^{-5} -19 \qr^{-7} \\
& \quad -7 \qr^{-9} +3 \qr^{-11} \\
[\lambda^{-8}]\colon\ & G^{4} \qr^{-5} +4 G^{3} \qr^{-4} +6 G^{3} \qr^{-6} +6 G^{2} \qr^{-3} +16 G^{2} \qr^{-5} +13 G^{2} \qr^{-7} \\
& \quad +4 G \qr^{-2} +14 G \qr^{-4} +20 G \qr^{-6} +12 G \qr^{-8} +\qr^{-1} +4 \qr^{-3} \\
& \quad +8 \qr^{-5} +9 \qr^{-7} +3 \qr^{-9} \\
[\lambda^{-6}]\colon\ & \qr^{-5} -\qr^{-7} \\
[\lambda^{-4}]\colon\ & \qr^{-3} -\qr^{-5} \\
[\lambda^{-2}]\colon\ & \qr^{-1} -\qr^{-3} \\
[\lambda^{0}]\colon\ & -\qr^{-1}
\end{align*}

$T(2,5)$、同じものを $\lambda^\beta$ ごとに因数分解した形（$\qr^2 - 1 = \bt^2 - 1$）
\begin{align*}
[\lambda^{-20}]\colon\ & \qr^{-17}\,\bigl(-1\bigr) \\
[\lambda^{-18}]\colon\ & \qr^{-17}\,\bigl(- (\qr^2 - 1)(G \qr + \qr^{2} + 2)\bigr) \\
[\lambda^{-16}]\colon\ & \qr^{-17}\,\bigl(- (\qr^2 - 1)(G \qr + \qr^{2} + 1)(G \qr^{3} + \qr^{4} + 2 \qr^{2} - 1)\bigr) \\
[\lambda^{-14}]\colon\ & \qr^{-15}\,\bigl(- (\qr^2 - 1) \\
& \quad (G^{3} \qr^{5} + 4 G^{2} \qr^{6} + 4 G^{2} \qr^{4} - G^{2} \qr^{2} \\
& \quad  + 5 G \qr^{7} + 10 G \qr^{5} + 2 G \qr^{3} - 3 G \qr + 2 \qr^{8} + 6 \qr^{6} \\
& \quad  + 5 \qr^{4} - 2 \qr^{2} - 1)\bigr) \\
[\lambda^{-12}]\colon\ & \text{既約（上の形のまま）} \\
[\lambda^{-10}]\colon\ & \text{既約（上の形のまま）} \\
[\lambda^{-8}]\colon\ & \text{既約（上の形のまま）} \\
[\lambda^{-6}]\colon\ & \qr^{-7}\,(\qr^2 - 1) \\
[\lambda^{-4}]\colon\ & \qr^{-5}\,(\qr^2 - 1) \\
[\lambda^{-2}]\colon\ & \qr^{-3}\,(\qr^2 - 1) \\
[\lambda^{0}]\colon\ & \qr^{-1}\,\bigl(-1\bigr)
\end{align*}


\section{unrefined の確認}\label{app:unref}
Mariño の trefoil の $\hat g$（標準 framing）を自前のコードで計算した値。$\hat g_{\yoneone}$ は \cite{Marino09} (4.8) の印刷値と全係数一致し、$\hat g_{\ytwo}$ は $z^9$ の項以外が一致する。
$\hat g_{\ytwo}$（$z = t_M - t_M^{-1}$ の冪ごと、$\lambda$ は Mariño の $\nu$）：
\begin{align*}
[z^{0}]\colon\ & -21 \lambda^{11} +79 \lambda^{9} -111 \lambda^{7} +69 \lambda^{5} -16 \lambda^{3} \\
[z]\colon\ & 21 \lambda^{12} -63 \lambda^{10} +63 \lambda^{8} -21 \lambda^{6} \\
[z^{2}]\colon\ & -70 \lambda^{11} +251 \lambda^{9} -307 \lambda^{7} +146 \lambda^{5} -20 \lambda^{3} \\
[z^{3}]\colon\ & 70 \lambda^{12} -231 \lambda^{10} +231 \lambda^{8} -70 \lambda^{6} \\
[z^{4}]\colon\ & -84 \lambda^{11} +330 \lambda^{9} -366 \lambda^{7} +128 \lambda^{5} -8 \lambda^{3} \\
[z^{5}]\colon\ & 84 \lambda^{12} -322 \lambda^{10} +322 \lambda^{8} -84 \lambda^{6} \\
[z^{6}]\colon\ & -45 \lambda^{11} +220 \lambda^{9} -230 \lambda^{7} +56 \lambda^{5} -\lambda^{3} \\
[z^{7}]\colon\ & 45 \lambda^{12} -219 \lambda^{10} +219 \lambda^{8} -45 \lambda^{6} \\
[z^{8}]\colon\ & -11 \lambda^{11} +78 \lambda^{9} -79 \lambda^{7} +12 \lambda^{5} \\
[z^{9}]\colon\ & 11 \lambda^{12} -78 \lambda^{10} +78 \lambda^{8} -11 \lambda^{6} \\
[z^{10}]\colon\ & -\lambda^{11} +14 \lambda^{9} -14 \lambda^{7} +\lambda^{5} \\
[z^{11}]\colon\ & \lambda^{12} -14 \lambda^{10} +14 \lambda^{8} -\lambda^{6} \\
[z^{12}]\colon\ & \lambda^{9} -\lambda^{7} \\
[z^{13}]\colon\ & -\lambda^{10} +\lambda^{8}
\end{align*}

$\hat g_{\yoneone}$（$z = t_M - t_M^{-1}$ の冪ごと、$\lambda$ は Mariño の $\nu$）：
\begin{align*}
[z^{0}]\colon\ & -15 \lambda^{11} +53 \lambda^{9} -69 \lambda^{7} +39 \lambda^{5} -8 \lambda^{3} \\
[z]\colon\ & 15 \lambda^{12} -45 \lambda^{10} +45 \lambda^{8} -15 \lambda^{6} \\
[z^{2}]\colon\ & -35 \lambda^{11} +126 \lambda^{9} -146 \lambda^{7} +61 \lambda^{5} -6 \lambda^{3} \\
[z^{3}]\colon\ & 35 \lambda^{12} -120 \lambda^{10} +120 \lambda^{8} -35 \lambda^{6} \\
[z^{4}]\colon\ & -28 \lambda^{11} +120 \lambda^{9} -128 \lambda^{7} +37 \lambda^{5} -\lambda^{3} \\
[z^{5}]\colon\ & 28 \lambda^{12} -119 \lambda^{10} +119 \lambda^{8} -28 \lambda^{6} \\
[z^{6}]\colon\ & -9 \lambda^{11} +55 \lambda^{9} -56 \lambda^{7} +10 \lambda^{5} \\
[z^{7}]\colon\ & 9 \lambda^{12} -55 \lambda^{10} +55 \lambda^{8} -9 \lambda^{6} \\
[z^{8}]\colon\ & -\lambda^{11} +12 \lambda^{9} -12 \lambda^{7} +\lambda^{5} \\
[z^{9}]\colon\ & \lambda^{12} -12 \lambda^{10} +12 \lambda^{8} -\lambda^{6} \\
[z^{10}]\colon\ & \lambda^{9} -\lambda^{7} \\
[z^{11}]\colon\ & -\lambda^{10} +\lambda^{8}
\end{align*}

